%! Minimizing Loss by Gradient Descent — Unit 8, Lesson 8.3 — Warm-Up
\documentclass[10pt]{article}
\usepackage{linalg-article}
\usepackage{linalg-boxes}
\newcommand{\TallMath}[1]{$\displaystyle #1\rule[-1.4em]{0pt}{3.2em}$}
\newcommand{\vv}{\mathbf{v}}

\begin{document}

\pageheader{Unit 8, Lesson 8.3}{Warm-Up}
\namedateperiod

\vspace{0.12in}

\begin{notesbox}{Getting Ready: Loss, Slope, and One Step}
\small Today we \emph{minimize the loss} --- roll a weight downhill until the error is smallest. That rests on
three moves: measuring a loss, reading a slope to find downhill, and updating with a step. Warm up on each.

\vspace{4pt}
\textbf{1. Squared-error loss (Unit 1 \& Lesson 8.0).} A network predicts $[3,5]$ but the targets are $[2,6]$.
The loss is the sum of squared errors:
\[
  L=(3-2)^2+(5-6)^2=\blank{1.4cm}.
\]
Smaller loss $=$ better fit. \emph{Learning} means making this number as small as possible.

\vspace{4pt}
\textbf{2. Which way is downhill? (slope).} A loss bowl has slope $L'(w)=2w-8$. Evaluate:
\[
  L'(2)=2(2)-8=\blank{1.0cm},\qquad L'(6)=2(6)-8=\blank{1.0cm}.
\]
A \emph{negative} slope means the loss still falls as $w$ grows (go right); a \emph{positive} slope means go left.

\vspace{4pt}
\textbf{3. Take one step (Unit 1 arithmetic).} The update rule is $w\leftarrow w-s\,L'(w)$. With current
$w=2$, learning rate $s=0.25$, and slope $L'(2)=-4$:
\[
  w = 2-0.25(-4)=\blank{1.2cm}.
\]
\end{notesbox}

\end{document}
